\textbf{Phần 1. Trắc nghiệm lựa chọn}
\setcounterpageref{ex}{0}

\begin{ex} %[1D9V34]
Đồng hồ có kim ngắn chỉ số 7, kim dài chỉ số 12. Hỏi đồng hồ chỉ mấy giờ?
\choice
{$6\text{ giờ}$}
{$8\text{ giờ}$}
{$5\text{ giờ}$}
{\True $7\text{ giờ}$}
\loigiai{
Kim dài chỉ số 12 và kim ngắn chỉ số 7: đồng hồ chỉ đúng 7 giờ.
\textbf{Đáp số:} $7$.\par
\textbf{Đáp án của câu hỏi là D}
}
\end{ex}

\begin{ex} %[1D9V34]
Đồng hồ có kim ngắn chỉ số 10, kim dài chỉ số 12. Hỏi đồng hồ chỉ mấy giờ?
\choice
{$0\text{ giờ}$}
{\True $10\text{ giờ}$}
{$20\text{ giờ}$}
{$9\text{ giờ}$}
\loigiai{
Đồng hồ chỉ đúng 10 giờ.
\textbf{Đáp số:} $10$.\par
\textbf{Đáp án của câu hỏi là B}
}
\end{ex}

\begin{ex} %[1D9V34]
Đồng hồ có kim ngắn chỉ số 3, kim dài chỉ số 12. Hỏi đồng hồ chỉ mấy giờ?
\choice
{$2\text{ giờ}$}
{$4\text{ giờ}$}
{\True $3\text{ giờ}$}
{$1\text{ giờ}$}
\loigiai{
Đồng hồ chỉ đúng 3 giờ.
\textbf{Đáp số:} $3$.\par
\textbf{Đáp án của câu hỏi là C}
}
\end{ex}

\begin{ex} %[1D9V34]
Đồng hồ có cả hai kim (kim ngắn và kim dài) cùng chỉ vào số 12. Hỏi đồng hồ chỉ mấy giờ?
\choice
{\True $12\text{ giờ}$}
{$11\text{ giờ}$}
{$13\text{ giờ}$}
{$10\text{ giờ}$}
\loigiai{
Cả hai kim cùng chỉ số 12: đồng hồ chỉ 12 giờ.
\textbf{Đáp số:} $12$.\par
\textbf{Đáp án của câu hỏi là A}
}
\end{ex}

\begin{ex} %[1D9V34]
Bé vào lớp học lúc 8 giờ sáng, tan học lúc 11 giờ trưa. Hỏi buổi sáng bé học ở trường trong bao nhiêu giờ?
\choice
{$2$}
{$4$}
{$1$}
{\True $3$}
\loigiai{
Thời gian bé học ở trường là: $11 - 8 = 3$ (giờ).
\textbf{Đáp số:} $3$.\par
\textbf{Đáp án của câu hỏi là D}
}
\end{ex}

\begin{ex} %[1D9V34]
Bố đi làm lúc 7 giờ sáng và về nhà lúc 12 giờ trưa. Hỏi bố đã đi làm trong bao nhiêu giờ?
\choice
{$4$}
{\True $5$}
{$6$}
{$3$}
\loigiai{
Thời gian bố đi làm là: $12 - 7 = 5$ (giờ).
\textbf{Đáp số:} $5$.\par
\textbf{Đáp án của câu hỏi là B}
}
\end{ex}

\begin{ex} %[1D9V34]
Buổi chiều bắt đầu từ 1 giờ chiều đến 5 giờ chiều. Khoảng thời gian này kéo dài bao nhiêu giờ?
\choice
{$3$}
{$5$}
{\True $4$}
{$2$}
\loigiai{
Khoảng thời gian này kéo dài: $5 - 1 = 4$ (giờ).
\textbf{Đáp số:} $4$.\par
\textbf{Đáp án của câu hỏi là C}
}
\end{ex}

\begin{ex} %[1D9V34]
Bạn Nam bắt đầu ngồi vào bàn học bài lúc 8 giờ tối và đi ngủ lúc 9 giờ tối. Nam học bài trong bao nhiêu giờ?
\choice
{\True $1$}
{$0$}
{$2$}
{$3$}
\loigiai{
Nam học bài trong: $9 - 8 = 1$ (giờ).
\textbf{Đáp số:} $1$.\par
\textbf{Đáp án của câu hỏi là A}
}
\end{ex}

\begin{ex} %[1D9V34]
Lúc 2 giờ chiều, sau 3 giờ nữa đồng hồ sẽ chỉ mấy giờ?
\choice
{$4\text{ giờ}$}
{$6\text{ giờ}$}
{$3\text{ giờ}$}
{\True $5\text{ giờ}$}
\loigiai{
Đồng hồ sẽ chỉ: $2 + 3 = 5$ (giờ chiều).
\textbf{Đáp số:} $5$.\par
\textbf{Đáp án của câu hỏi là D}
}
\end{ex}

\begin{ex} %[1D9V34]
Lúc 9 giờ sáng, trước đó 2 giờ là mấy giờ?
\choice
{$6\text{ giờ}$}
{\True $7\text{ giờ}$}
{$8\text{ giờ}$}
{$5\text{ giờ}$}
\loigiai{
Thời điểm trước đó 2 giờ là: $9 - 2 = 7$ (giờ sáng).
\textbf{Đáp số:} $7$.\par
\textbf{Đáp án của câu hỏi là B}
}
\end{ex}

\begin{ex} %[1D9V34]
Mặt đồng hồ có các số từ 1 đến bao nhiêu?
\choice
{$11$}
{$13$}
{\True $12$}
{$10$}
\loigiai{
Mặt đồng hồ tròn có các số từ 1 đến 12.
\textbf{Đáp số:} $12$.\par
\textbf{Đáp án của câu hỏi là C}
}
\end{ex}

\begin{ex} %[1D9V34]
Một ngày có bao nhiêu giờ?
\choice
{\True $24$}
{$23$}
{$25$}
{$22$}
\loigiai{
Một ngày đêm có tất cả $24$ giờ.
\textbf{Đáp số:} $24$.\par
\textbf{Đáp án của câu hỏi là A}
}
\end{ex}

\textbf{Phần 2. Trắc nghiệm đúng sai}
\setcounterpageref{ex}{0}

\begin{ex} %[1D9V34]
Xét tính đúng hoặc sai của các khẳng định toán học sau đây:
\choiceTF
{Kết quả của bài toán «Kim dài trên mặt đồng hồ dùng để chỉ gì? (1: Giờ; 2: Phút; 3: Giây)» là $3\text{ giờ}$.}
{Kết quả của bài toán «Kim ngắn trên mặt đồng hồ dùng để chỉ gì? (1: Giờ; 2: Phút; 3: Giây)» là $0\text{ giờ}$.}
{\True Kết quả của bài toán «Khi đồng hồ chỉ giờ đúng, kim dài luôn luôn chỉ vào số nào» là $12\text{ giờ}$.}
{\True Kết quả của bài toán «Xe buýt khởi hành lúc 6 giờ sáng và đến nơi lúc 8 giờ sáng. Chuyến xe buýt chạy hết bao nhiêu giờ» là $2$.}
\loigiai{
\begin{itemize}
    \item Ý a) Sai: Kết quả đúng là $2$, không phải $3$. Kim dài là kim chỉ phút (ghi số 2).
\textbf{Đáp số:} $2$.
    \item Ý b) Sai: Kết quả đúng là $1$, không phải $0$. Kim ngắn là kim chỉ giờ (ghi số 1).
\textbf{Đáp số:} $1$.
    \item Ý c) Đúng: Khi đồng hồ chỉ giờ đúng, kim dài luôn chỉ vào số $12$.
\textbf{Đáp số:} $12$.
    \item Ý d) Đúng: Thời gian xe buýt chạy là: $8 - 6 = 2$ (giờ).
\textbf{Đáp số:} $2$.
\end{itemize}
\begin{itemize}
    \item a) \textbf{Sai}.
    \item b) \textbf{Sai}.
    \item c) \textbf{Đúng}.
    \item d) \textbf{Đúng}.
\end{itemize}
}
\end{ex}

\begin{ex} %[1D9V34]
Xét tính đúng hoặc sai của các bài toán và phép tính sau đây:
\choiceTF
{\True Mẹ nấu cơm mất 1 giờ. Mẹ bắt đầu nấu lúc 10 giờ trưa. cơm chín lúc $11\text{ giờ}$ giờ.}
{Kết quả của bài toán «Gia đình bạn An ăn cơm tối lúc 7 giờ tối. Sau 2 giờ, cả nhà cùng xem phim. Cả nhà xem phim lúc mấy giờ» là $8\text{ giờ}$.}
{Bộ phim hoạt hình dài 1 giờ, kết thúc lúc 5 giờ chiều. bộ phim bắt đầu chiếu lúc $5\text{ giờ}$ giờ.}
{\True Kết quả của bài toán «Khi đồng hồ chỉ giờ đúng, kim dài luôn luôn chỉ vào số nào» là $12\text{ giờ}$.}
\loigiai{
\begin{itemize}
    \item Ý a) Đúng: Cơm chín lúc: $10 + 1 = 11$ (giờ trưa).
\textbf{Đáp số:} $11$.
    \item Ý b) Sai: Kết quả đúng là $9$, không phải $8$. Cả nhà xem phim lúc: $7 + 2 = 9$ (giờ tối).
\textbf{Đáp số:} $9$.
    \item Ý c) Sai: Kết quả đúng là $4$, không phải $5$. Bộ phim bắt đầu lúc: $5 - 1 = 4$ (giờ chiều).
\textbf{Đáp số:} $4$.
    \item Ý d) Đúng: Khi đồng hồ chỉ giờ đúng, kim dài luôn chỉ vào số $12$.
\textbf{Đáp số:} $12$.
\end{itemize}
\begin{itemize}
    \item a) \textbf{Đúng}.
    \item b) \textbf{Sai}.
    \item c) \textbf{Sai}.
    \item d) \textbf{Đúng}.
\end{itemize}
}
\end{ex}

\textbf{Phần 3. Trả lời ngắn}
\setcounterpageref{ex}{0}

\begin{ex} %[1D9V34]
Đồng hồ báo thức reo lúc 6 giờ sáng. Bạn Minh dậy tập thể dục trong 1 giờ. Minh tập thể dục xong lúc mấy giờ?
\shortans{7}
\loigiai{
\textbf{Đáp số:} $7$\par
Minh tập thể dục xong lúc: $6 + 1 = 7$ (giờ sáng).
\textbf{Đáp số:} $7$.
}
\end{ex}

\begin{ex} %[1D9V34]
Đồng hồ A chỉ 3 giờ, đồng hồ B chỉ 6 giờ. Hỏi đồng hồ B chỉ muộn hơn đồng hồ A bao nhiêu giờ?
\shortans{3}
\loigiai{
\textbf{Đáp số:} $3$\par
Đồng hồ B chỉ muộn hơn đồng hồ A là: $6 - 3 = 3$ (giờ).
\textbf{Đáp số:} $3$.
}
\end{ex}

\begin{ex} %[1D9V34]
Đồng hồ chỉ 9 giờ. Kim ngắn và kim dài tạo với nhau thành một góc gì? (1: Góc nhọn; 2: Góc tù; 3: Góc vuông; 4: Góc bẹt).
\shortans{3}
\loigiai{
\textbf{Đáp số:} $3$\par
Lúc 9 giờ đúng, kim ngắn chỉ số 9 và kim dài chỉ số 12 tạo thành góc vuông (ghi số 3).
\textbf{Đáp số:} $3$.
}
\end{ex}

\begin{ex} %[1D9V34]
Đồng hồ chỉ 3 giờ. Hai kim đồng hồ tạo với nhau thành một góc gì? (1: Góc nhọn; 2: Góc tù; 3: Góc vuông; 4: Góc bẹt).
\shortans{3}
\loigiai{
\textbf{Đáp số:} $3$\par
Lúc 3 giờ đúng, hai kim tạo thành góc vuông (ghi số 3).
\textbf{Đáp số:} $3$.
}
\end{ex}

\begin{ex} %[1D9V34]
Đồng hồ chỉ 6 giờ. Hai kim đồng hồ thẳng hàng tạo thành một góc gì? (1: Góc nhọn; 2: Góc tù; 3: Góc vuông; 4: Góc bẹt).
\shortans{4}
\loigiai{
\textbf{Đáp số:} $4$\par
Lúc 6 giờ đúng, kim ngắn chỉ số 6, kim dài chỉ số 12 tạo thành đường thẳng (góc bẹt, ghi số 4).
\textbf{Đáp số:} $4$.
}
\end{ex}

\begin{ex} %[1D9V34]
Từ 8 giờ sáng đến 10 giờ sáng, kim giờ đã đi được bao nhiêu khoảng số trên mặt đồng hồ?
\shortans{2}
\loigiai{
\textbf{Đáp số:} $2$\par
Kim giờ đi từ số 8 đến số 10 được: $10 - 8 = 2$ (khoảng số).
\textbf{Đáp số:} $2$.
}
\end{ex}
